%%─────────────────────────────────────────────────────────────────────────────
\chapter{USA}
%%─────────────────────────────────────────────────────────────────────────────

\section{The Constitutional Foundation}



The 10th Amendment is kind of The Root of Everything. The US Constitution (ratified 1788) does \textbf{not} mention education anywhere \cite{us-constitution}.

\begin{quote}
	``The powers not delegated to the United States by the Constitution, nor prohibited by it to the States, are reserved to the States respectively, or to the people.'' \cite{us-constitution}
\end{quote}

This one sentence is why the USA has \textbf{no national curriculum, no national textbook board, and no national textbook.} \cite{us-constitution} Education is constitutionally a \textbf{state power}, not a federal power.

\vspace{0.5cm}



The USA has \textbf{50 states}. Each state legally operates its own education system \cite{us-ed-gov-k12}. Technically, there is no single ``American education system''. There are \textbf{50 systems} running in parallel, with some federal overlay (explained in Section 2). The US Department of Education (DoED) was created by Congress in 1979 via the Department of Education Organization Act \cite{us-doed-act}. It does \textbf{not} write curriculum, mandate textbooks, set graduation requirements, or run any school. The Federal Government distributes federal \textbf{funding} to states (via formulas and grants) \cite{us-ed-gov-k12}, enforces federal \textbf{civil rights laws} in schools (e.g., no discrimination) \cite{us-doed-act}, collects education data and statistics nationally \cite{nces-digest}, runs federal student loan programs for higher education \cite{nces-ipeds}, and sets conditions tied to federal money (see Section 3). Under the Trump administration (returned January 2025), the DoED has been significantly restructured \cite{us-doed-2026}. An executive order was signed targeting its dismantlement. Many functions were transferred to other departments (e.g., student loans to Treasury, special education to Health and Human Services). As of mid-2026, the DoED exists in a drastically reduced form. \textbf{This is an active, ongoing situation --- not yet fully resolved by legislation.} Even when the DoED existed at full power, it never had the authority. Its dismantlement makes curriculum even more decentralized. Each of the 50 states has a State Board of Education and a State Department of Education \cite{ref18}. These bodies write or adopt academic content standards, set graduation requirements, set teacher certification requirements, determine school year length, and handle state funding distribution to local districts.

\vspace{0.3cm}

Academic Standards are \textbf{not} textbooks \cite{us-ed-gov-k12}. They are documents that say, for example: ``By the end of 4th grade, students should be able to multiply two-digit numbers.'' Standards tell you \textbf{what} to teach. They do not tell you \textbf{which book} to use to teach it.

\vspace{0.3cm}
{\large\textbf{The Common Core Experiment (2010--present)}}

In 2010, the National Governors Association and the Council of Chief State School Officers (non-government bodies) created the Common Core State Standards (CCSS) for Math and English Language Arts. These standards were voluntary—no law forced states to adopt them, though the federal government incentivized adoption via the Race to the Top grant program. At peak adoption, 45 states plus DC adopted Common Core. However, as of 2023–2024, several states withdrew (including Indiana in 2014, Oklahoma in 2014, and South Carolina in 2015, while Texas, Alaska, and Virginia never adopted them). Ultimately, Common Core is not a national curriculum, but rather a voluntary alignment attempt that partially fragmented.

Research shows: in approximately \textbf{43 states}, curriculum selection authority rests primarily with \textbf{local school districts}, not the state government \cite{textbook-adoption-states}. This means: even within one state, different school districts can teach using different textbooks and curricula.

%%─────────────────────────────────────────────────────────────────────────────

A school district is a local educational agency (LEA) — a government body that directly runs public schools in a defined geographic area. The USA has approximately 13,000 school districts (US Census Bureau data). Each district has a locally elected or appointed School Board that decides which textbooks to purchase, which curriculum programs to implement, and local school policies and tax levies. In districts with local curriculum authority, a 7th-grade math teacher in one Texas district might use a completely different textbook than a teacher in a neighboring district \cite{textbook-adoption-states}. There is no law against this.

\textbf{Adoption States vs.\ Open Territory States}

\textbf{Adoption States ($\sim$19--20 states):} The state government maintains an \textbf{official approved list} of textbooks for K-12 subjects \cite{textbook-adoption-states}. Districts must select from that approved list. Examples: Texas, California, Florida, North Carolina, Virginia.

\textbf{Open Territory States ($\sim$30--31 states):} No state-approved list \cite{textbook-adoption-states}. Local districts choose any textbook from any publisher.

Texas ($\sim$5.4M K-12 students) and California ($\sim$6M K-12 students) are the two largest state markets \cite{nces-digest}. Because they are adoption states with large budgets, textbook publishers \textbf{design national editions to satisfy Texas and California first}. They then sell the same books to the rest of the country. Texas and California effectively shape what's in American textbooks nationwide, even though they have no legal authority over other states \cite{textbook-adoption-states}. The USA has a \textbf{private, for-profit textbook publishing industry} \cite{textbook-adoption-states}. Major publishers: Pearson, McGraw-Hill, Houghton Mifflin Harcourt (HMH), Bedford Freeman \& Worth, Cengage. These are corporations. Schools and districts \textbf{purchase} textbooks with public (tax) funds. In public K-12, students use textbooks owned by the school and return them at year end \cite{us-ed-gov-k12}. Students do not usually buy or keep them.

The K-12 Grade Structure consists of Kindergarten for grade K with approximate ages of 5 through 6, Elementary School for grades 1 through 5 or 1 through 6 with approximate ages of 6 through 11, Middle School or Junior High for grades 6 through 8 or 7 through 8 with approximate ages of 11 through 14, and High School for grades 9 through 12 with approximate ages of 14 through 18. ``K-12'' = Kindergarten through Grade 12. Approximately 13 years of schooling before university \cite{us-ed-gov-k12}. The exact split varies by state and district \cite{us-ed-gov-k12}. Some places run K-8 then 9-12 high school. There is no single mandated structure nationally. Each state may have its own \textbf{state assessment tests} (STAAR in Texas, MCAS in Massachusetts, PARCC in some states). These vary enormously. Historically, students took the \textbf{SAT} (College Board) or \textbf{ACT} (ACT Inc.) for college admissions --- but these are \textbf{private standardized tests}, not government exams, and many universities have now made them optional \cite{sat-act}.

\begin{center}\small
	\begin{tabularx}{0.97\textwidth}{@{} >{\raggedright\arraybackslash}p{3.5cm} X @{}}
		\toprule
		\textbf{Type}         & \textbf{Description}                                                  \\
		\midrule
		Community College     & 2-year public colleges, offer Associate degrees, low cost             \\
		Liberal Arts College  & 4-year private colleges, small, focus on undergraduate education      \\
		State University      & Public 4-year universities, funded by state government                \\
		Research University   & Public or private, offer Bachelor + Master + PhD, major research      \\
		For-Profit University & Private companies operating universities (controversial, many closed) \\
		\bottomrule
	\end{tabularx}
\end{center}

The USA has approximately \textbf{5,000+ degree-granting institutions} \cite{nces-ipeds}. Each American university grants its \textbf{own degree} in its own name. No central body grants the actual degree. Each university runs its own \textbf{individual admissions process} evaluating: high school GPA and transcript, SAT/ACT scores (if required), essays, extracurriculars, and letters of recommendation. Thousands of universities = thousands of separate applications and processes. Universities choose their own degree programs, textbooks (professors choose individually), and assessment methods \cite{nces-ipeds}. A professor teaching Introductory Statistics at University of Michigan can assign a completely different textbook than a professor at UCLA. \textbf{University textbooks are expensive private market products.} A single textbook can cost \$100--\$300+ USD \cite{textbook-cost-crisis}. This is a well-documented crisis in American higher education.

Since the government does not run universities, quality is controlled through \textbf{accreditation} --- a peer-review process conducted by private non-profit bodies \cite{accreditation-chea}. The Seven Regional Accrediting Bodies include NECHE for New England, MSCHE for Middle States, HLC for the large North Central area, SACSCOC for Southern states, WSCUC for the Western region including California, Hawaii, and the Pacific, NWCCU for the Northwest, and ACCJC for California community colleges.
A university \textbf{must} be regionally accredited for its degrees to be widely recognized and for students to use federal financial aid \cite{accreditation-chea}. Unaccredited institutions' degrees are generally not recognized by employers or other universities. These bodies are \textbf{private non-profits}, not government agencies \cite{accreditation-chea}. The government recognizes them but does not operate them. This is a fundamentally American model --- private self-regulation with government recognition.



\section{The 10th Amendment --- The Root of Everything}
\section{There Are 50 Separate Education Systems}
\section{The Federal Government's Role}
\subsection{The US Department of Education (DoED)}
\subsection{What the Federal Government Actually Does}
\subsection{Current Status of the DoED (July 2026)}

\section{State-Level Authority}
\subsection{States Set the Standards}
\subsection{Academic Standards --- What Are They?}
\subsection{The Common Core Experiment (2010--present)}
\subsection{State Curriculum Authority}
\subsection{Every Student Succeeds Act (ESSA 2015)}

\section{Local School Districts}
\subsection{What Is a School District?}
\subsection{Implication: Textbook Chaos (Or Freedom)}

\section{Public vs.\ Charter vs.\ Private Schools}
\subsection{Public Schools}
\subsection{Charter Schools}
\subsection{Private Schools}
\subsection{Homeschooling}

\section{The Textbook System in America}
\subsection{Adoption States vs.\ Open Territory States}
\subsection{Why Texas and California Are Disproportionately Powerful}
\subsection{The Private Textbook Publishing Industry}
\subsection{No Free Textbooks in Most of K-12}

\section{The K-12 Grade Structure}
\subsection{Basic Structure}
\subsection{Variation Exists}
\subsection{No National School Leaving Exam}
\subsection{SAT and ACT --- Private Standardized Tests}

\section{Special Education}
\subsection{Individuals with Disabilities Education Act (IDEA)}
\subsection{Inclusion Model}

\section{American Higher Education}
\subsection{Structure of Higher Education}
\subsection{No National University}
\subsection{No National University Entrance Exam}
\subsection{No National Higher Education Curriculum}
\subsection{The Textbook Cost Crisis}
\subsection{Student Loan System}

\section{Accreditation}
\subsection{What Replaces Government Certification?}
\subsection{The Seven Regional Accrediting Bodies}
\subsection{Accreditation Is Not Government-Run}


%%─────────────────────────────────────────────────────────────────────────────
\chapter{Japan}
%%─────────────────────────────────────────────────────────────────────────────

\section{Legal and Constitutional Basis}
\subsection{Ministry of Education, Culture, Sports, Science and Technology (MEXT)}
\subsection{Fundamental Law of Education (2006 Revision)}

\section{K-12 School Structure}
\subsection{Elementary School --- Shogakko (Grades 1--6)}
\subsection{Middle School --- Chugakko (Grades 7--9)}
\subsection{High School --- Kotogakko (Grades 10--12)}
\subsection{Compulsory Education Scope}

\section{National Curriculum Standards}
\subsection{Course of Study (Gakushu Shido Yoryo)}
\subsection{How Often It Is Revised}
\subsection{Textbook Authorization System}
\subsection{Textbook Distribution (Free for Compulsory Levels)}

\section{Exam Culture}
\subsection{High School Entrance Exams}
\subsection{University Entrance Exam --- NCEA (Daigaku Nyushi Kyotsu Test)}
\subsection{Individual University Exams}
\subsection{Juku --- Cram Schools}

\section{Character Education}
\subsection{Tokkatsu --- Classroom Activities Beyond Academics}
\subsection{Moral Education (Dotoku)}
\subsection{Club Activities (Bukatsu)}

\section{Teacher System}
\subsection{Teacher Certification}
\subsection{Teacher Autonomy vs.\ Curriculum Control}
\subsection{Teacher Training}

\section{Special and Vocational Education}
\subsection{Special Needs Education}
\subsection{Vocational High Schools}
\subsection{Vocational Colleges --- Senmon Gakko}

\section{Higher Education}
\subsection{Types of Institutions}
\subsection{National vs.\ Public vs.\ Private Universities}
\subsection{Graduate Education}
\subsection{Tuition and Funding}

\section{English Education Reform (2020--present)}
\section{International Baccalaureate Adoption in Japan}
\section{Current Challenges (As of 2026)}


%%─────────────────────────────────────────────────────────────────────────────
\chapter{Finland}
%%─────────────────────────────────────────────────────────────────────────────

\section{Legal and Constitutional Basis}
\subsection{Finnish National Agency for Education (EDUFI/OPH)}
\subsection{Ministry of Education and Culture}

\section{Core Philosophy}
\subsection{Equity Over Excellence}
\subsection{Trust Over Accountability}
\subsection{Wellbeing Over Performance}

\section{K-12 School Structure}
\subsection{Early Childhood Education and Care (ECEC)}
\subsection{Pre-Primary Education (Age 6)}
\subsection{Primary School --- Alakoulu (Grades 1--6)}
\subsection{Lower Secondary --- Yläkoulu (Grades 7--9)}
\subsection{Upper Secondary --- Lukio (Grades 10--12)}
\subsection{Vocational Upper Secondary Education}

\section{National Core Curriculum}
\subsection{How It Works}
\subsection{Local Curriculum Autonomy}
\subsection{Phenomenon-Based Learning}

\section{Assessment Philosophy}
\subsection{No Standardized Testing in Grades 1--9}
\subsection{Teacher Assessment Model}
\subsection{Matriculation Examination --- Ylioppilastutkinto}

\section{Teacher System}
\subsection{Teacher as Profession: Master's Degree Required}
\subsection{Teacher Autonomy in Material Selection}
\subsection{Teacher Training and Selection}

\section{Textbook System}
\subsection{No State-Mandated Textbooks}
\subsection{Private Publishers, Teacher Choice}

\section{Special Education and Inclusion}
\subsection{Three-Tier Support Model}
\subsection{Early Intervention Philosophy}

\section{Multilingual and Immigrant Education}
\section{School Funding Model}

\section{Higher Education}
\subsection{Universities --- Research Focus}
\subsection{Universities of Applied Sciences --- Ammattikorkeakoulut (AMK)}
\subsection{Tuition-Free Education}
\subsection{Student Financial Aid (KELA)}

\section{PISA Rankings --- Context and Criticism}
\section{Current Challenges (As of 2026)}


%%─────────────────────────────────────────────────────────────────────────────
\chapter{Netherlands and the Bologna Process}
%%─────────────────────────────────────────────────────────────────────────────

\section{The Bologna Process}
\subsection{What Is It?}
\subsection{European Higher Education Area (EHEA)}
\subsection{The Three-Cycle Degree System}
\subsection{ECTS Credit System}
\subsection{What Bologna Does NOT Do}
\subsection{49 Member Countries (As of 2024)}

\section{Legal and Constitutional Basis}
\subsection{Freedom of Education --- Article 23 of Dutch Constitution}
\subsection{Ministry of Education, Culture and Science (OCW)}

\section{The Pillarization System (Verzuiling)}
\subsection{Public vs.\ Religiously Affiliated Schools}
\subsection{Equal Funding Regardless of Type}

\section{K-12 School Structure}
\subsection{Primary School --- Basisonderwijs (Ages 4--12)}
\subsection{Secondary Education --- Voortgezet Onderwijs (VO)}
\subsubsection{VMBO --- Vocational Preparation (Ages 12--16)}
\subsubsection{HAVO --- Higher General (Ages 12--17)}
\subsubsection{VWO --- Pre-University (Ages 12--18)}
\subsection{Central Final Exams (Eindexamen)}

\section{Curriculum and School Autonomy}
\subsection{Core Objectives Set by Government}
\subsection{Schools Choose Their Own Methods and Materials}
\subsection{No National Textbook List}

\section{Teacher System}
\subsection{Teacher Certification Requirements}
\subsection{Teacher Autonomy}

\section{Special Education}
\subsection{Appropriate Education Act (Passend Onderwijs)}

\section{Higher Education}
\subsection{Research Universities --- Universiteiten (WO)}
\subsection{Universities of Applied Sciences --- Hogescholen (HBO)}
\subsection{Tuition Fees and Student Funding}
\subsection{Internationalization}
\subsection{Quality Assurance and Accreditation (NVAO)}

\section{Student Mobility Within EHEA}
\section{Current Challenges (As of 2026)}


%%─────────────────────────────────────────────────────────────────────────────
\chapter{Key Takeaways}
%%─────────────────────────────────────────────────────────────────────────────

\section{Decentralization vs.\ Centralization Spectrum}
\section{Testing and Assessment Philosophies}
\section{Teacher Autonomy Spectrum}
\section{Equity vs.\ Excellence}
\section{Textbook Systems Compared}
\section{Higher Education Access and Cost}
\section{Bologna Process Impact}
\section{What Bangladesh Can Learn}
